% ============================================================
%  Chapitre 2 — Modèle et implémentation de référence
% ============================================================
\chapter{Modèle ACO et implémentation de référence}
\label{chap:modele}

\section{Description formelle du modèle}

\subsection{Représentation de l'environnement}

L'environnement est une grille cartésienne $N \times N$ ($N = 512$). Chaque cellule $s$
stocke trois valeurs :
\begin{itemize}
    \item $c(s) \in [0, 1]$ : coût de traversée (généré par plasma fractal) ;
    \item $V_1(s) \in [0, 1]$ : phéromone de type 1, guide vers la \emph{nourriture} ;
    \item $V_2(s) \in [0, 1]$ : phéromone de type 2, guide vers le \emph{nid}.
\end{itemize}

\subsection{Dynamique des phéromones}

À chaque pas de temps, lorsqu'une fourmi visite la cellule $s$, les phéromones sont mis à jour
selon :

\begin{equation}
    V_i(s) \leftarrow
    \begin{cases}
        1 & \text{si } s \text{ est la source cible du phéromone } i \\
        \alpha \cdot \max_{s' \in \mathcal{N}(s)} V_i(s')
        + (1 - \alpha) \cdot \frac{1}{4}\sum_{s' \in \mathcal{N}(s)} V_i(s')
        & \text{sinon}
    \end{cases}
    \label{eq:pheromone_update}
\end{equation}

où $\mathcal{N}(s)$ désigne les quatre voisins (haut, bas, gauche, droite) de $s$,
et $\alpha \in [0,1]$ est le \textbf{paramètre de bruit}.

Après chaque itération, tous les phéromones s'évaporent :
\begin{equation}
    V_i(s) \leftarrow \beta \cdot V_i(s), \quad \forall s
    \label{eq:evaporation}
\end{equation}
où $\beta \in ]0, 1[$ est le \textbf{coefficient d'évaporation}.

\subsection{Règle de déplacement}

À chaque sous-pas, une fourmi en position $s$ choisit sa prochaine cellule selon :
\begin{itemize}
    \item avec probabilité $\varepsilon$ (\textbf{exploration}) : cellule voisine tirée
          uniformément parmi $\mathcal{N}(s) \setminus \{\text{obstacles}\}$ ;
    \item avec probabilité $1 - \varepsilon$ (\textbf{exploitation}) : cellule $s^*$ maximisant
          $V_{i^*}(s')$, où $i^* = 1$ si la fourmi est non chargée, $i^* = 2$ si elle est chargée.
\end{itemize}

La fourmi avance tant que son budget de mouvement $b$ (initialisé à 1) n'est pas épuisé :
$b \leftarrow b - c(s_{\text{new}})$.

\subsection{Paramètres}

\begin{table}[h]
\centering
\caption{Paramètres de la simulation de référence}
\label{tab:params}
\begin{tabular}{llll}
\hline
\textbf{Paramètre} & \textbf{Symbole} & \textbf{Valeur} & \textbf{Rôle} \\
\hline
Nombre de fourmis    & $m$          & 5\,000  & Agents simulés \\
Taille de la grille  & $N$          & 512     & $N \times N$ cellules \\
Exploration          & $\varepsilon$& 0.8     & Proportion d'exploration \\
Bruit                & $\alpha$     & 0.7     & Pondération max/moyenne \\
Évaporation          & $\beta$      & 0.999   & Décroissance des phéromones \\
\hline
\end{tabular}
\end{table}

\section{Génération du terrain fractal}

Le terrain est généré par l'algorithme \textit{midpoint displacement} récursif :
une grille de $n \times n$ sous-grilles de taille $2^k + 1$ est initialisée aux coins,
puis subdivisée récursivement en calculant les points médians avec une déviation bornée par
$d \times \frac{ns}{2}$. Les valeurs sont ensuite normalisées dans $[0, 1]$.

% TODO : ajouter une figure du terrain fractal généré (capture d'écran SDL)
% \begin{figure}[h]
%     \centering
%     \includegraphics[width=0.6\textwidth]{Figures/fractal_land.png}
%     \caption{Exemple de terrain fractal 512×512 (seed=1024, deviation=1.0)}
%     \label{fig:fractal}
% \end{figure}

\section{Architecture du code de référence}

Le code est organisé en quatre modules :

\begin{description}
    \item[\texttt{fractal\_land}] Génération et stockage du terrain (tableau 1D de \texttt{double}).
    \item[\texttt{pheronome}] Grille $(N+2)^2$ de paires $[V_1, V_2]$, avec cellules fantômes
          à $-1$ pour les bords. Méthodes : \texttt{mark\_pheronome()}, \texttt{do\_evaporation()},
          \texttt{update()} (double-buffering).
    \item[\texttt{ant}] Un objet par fourmi (\textbf{AoS}) : état, position (\texttt{SDL\_Point}),
          graine aléatoire. Méthode \texttt{advance()} : boucle de sous-pas jusqu'à budget épuisé.
    \item[\texttt{ant\_simu}] Point d'entrée, boucle principale SDL, appel à
          \texttt{advance\_time()} à chaque itération.
\end{description}

\subsection{Analyse de complexité}

La complexité par itération est dominée par la phase de déplacement des fourmis :
$\mathcal{O}(m \cdot \bar{k})$ où $\bar{k}$ est le nombre moyen de sous-pas par fourmi
(dépend du terrain et de la position). La phase d'évaporation est $\mathcal{O}(N^2)$, fixe
et régulière.

% TODO : compléter avec une mesure expérimentale de k_bar moyen
